\chapter{Literature Review}

\section{Introduction to Routing Algorithms in NoCs}

Routing determines how packets travel through a Network-on-Chip from their source to their destination. In a typical 2D mesh NoC, each router forwards packets according to the routing algorithm implemented in its routing unit, with the overall goal of minimizing latency, sustaining throughput, and avoiding congestion and deadlock. NoC routing algorithms are broadly categorized as static or adaptive. Static routing uses fixed, predetermined paths independent of runtime conditions. Deterministic XY routing is the standard example, where packets are routed first in the X direction and then in the Y direction. Although this approach is simple and low cost, it does not react to traffic imbalance or failures, and can therefore create bottlenecks and underutilize alternative paths \cite{AROMa2018}.

Adaptive routing addresses this limitation by selecting routes based on the current network state, such as congestion levels, fault status, or buffer availability. These routing decisions may rely on local information, global information, or a hybrid of both, depending on the algorithm design. For this work, we are interested in adaptive routing only, which is best understood through the taxonomy proposed in the recent survey of 2D NoC routing \cite{Kondoth2026}, which divides adaptive approaches into reactive, proactive, and application-specific classes. Reactive routing responds to conditions that already exist in the network, while proactive routing attempts to prevent future degradation by anticipating harmful operating conditions.

Within reactive adaptation, one important class is fault-tolerant routing, whose objective is to preserve packet delivery by rerouting traffic around failed routers or links. Within proactive adaptation, one important class is aging-aware routing, which aims to reduce long-term wear by steering traffic away from components experiencing higher stress. In other words, fault-tolerant routing responds after faults occur, whereas aging-aware routing attempts to delay those faults by controlling the accumulation of wear over time \cite{Kondoth2026}. This work focuses on extending NoC lifetime as opposed to merely recovering from failure. The remainder of this literature review concentrates on adaptive routing methods in these two categories, with particular emphasis on aging-aware routing.

\section{Fault-Tolerant Routing Algorithm in NoC}

Fault-tolerant routing is concerned with maintaining communication in the presence of faulty routers or links. Its primary objective is to minimize packet loss and maximize successful packet delivery by rerouting traffic around failed components, while keeping latency and throughput within acceptable limits.

A simple approach is the dynamic XY-YX fault-tolerant routing algorithm \cite{An2021}, which switches from XY routing to YX routing when faults are encountered. However, the method is fundamentally constrained by a small set of predefined routing responses. Once faults create more complex blockage patterns, the algorithm cannot reason beyond the fixed XY/YX alternatives. Similar approaches are represented by logic- and rule-based fault-tolerant routing schemes summarized in the survey \cite{Kondoth2026}, such as FTMR, RAFT, Hermes, and related fault-ring or contour-based methods. These approaches identify faulty regions, classify unusable nodes, or maintain local fault information to bypass failed components. The drawback of these methods is that they depend on explicit fault modeling, packet-header updates, or region management logic, all of which add control overhead and can still lead to suboptimal paths when network conditions become more complex.

Recent work has therefore explored reinforcement learning (RL) as a way to make fault-tolerant routing more adaptive. In \cite{Jagadheesh2020}, routing is modeled as a Markov Decision Process and solved using Q-learning, with routers treated as states and routing decisions as actions. The proposed method is shown to find optimal paths in the presence of both router and link faults. On average, the paper reports about 19\% more delivered packets under link faults and 11\% more delivered packets under router faults. However, it relies on a large Q-table that grows with the number of routers and possible fault conditions, which can lead to high memory overhead. This limitation is addressed more directly in \cite{Samala2020}. This work extends the earlier RL-based mesh routing approach using decentralized multi-agent reinforcement learning and a two-level Q-table. The goal is to reduce the storage cost and stale-value issues associated with large flat Q-tables, while still learning effective fault-tolerant policies. More recent work \cite{LARE2023} shows that RL can improve network performance by learning better flit rerouting behavior.

\section{Aging-Aware Routing Algorithm}

Aging-aware routing aims to extend NoC lifetime by steering traffic away from components that are aging faster due to stress, temperature, and process variation. Unlike fault-tolerant routing, which reacts after failures occur, aging-aware routing attempts to delay failure by reducing the uneven accumulation of wear across routers and links. As the survey notes \cite{Kondoth2026}, the main evaluation metric in this class of work is typically MTTF, since the central objective is long-term reliability rather than immediate fault recovery.

One of the most influential works in this area is Clotho \cite{Clotho2015}, which frames wear-aware routing as a path-selection problem based on the lifetime of candidate routes. Clotho models both EM and HCI, while also incorporating temperature and process variation. Using these inputs, it computes a path metric over minimal routes and directs traffic toward paths with higher relative lifetime. Under PARSEC workloads, Clotho reports average MTTF improvements of 66\% for EM and 8\% for HCI over deterministic routing. However, Clotho relies on explicit path evaluation and "piggybacking" network lifetime information across the network, which introduces computational and memory overhead, and limits scalability.

A different line of work focuses on combining aging mitigation with performance control. Bhardwaj et al. \cite{Bhardwaj2012} introduces the Traffic Threshold per Epoch (TTpE) metric, which limits how much traffic a stressed router or link should accept during a given runtime interval. Its routing algorithm chooses paths that are both less stressed and less congested, and can insert recovery cycles once a component reaches its threshold. Experimentally, the paper reports reductions of 13\% in network latency, 12.17\% in EDPPF, and a 10.4\% performance improvement over its comparison scheme. However, this method relies on precomputed lifetime models and static thresholds, which may not adapt well to changing runtime conditions or workloads.

Another important direction is online monitoring-based aging awareness. In \cite{Ghaderi2017}, Ghaderi et al. propose a Centralized Aging Table (CAT) that maps observed router activity to aging degradation. Their method monitors the number of flits and residence time within routers, then periodically updates routing to avoid highly aged routers by selecting among shortest paths. Using gem5, it reports average improvements of 39\% in maximum-aged router degradation and 52\% in aging imbalance compared to XY routing, while introducing negligible overhead in latency and EDP. This idea is extended further in AROMa \cite{AROMa2018}, which targets 3D NoCs. It uses a Distributed Centralized Aging Table (D-CAT), with a separate table per layer, and periodically swaps among k-best shortest paths to avoid highly aged routers and allow BTI recovery. Compared to non-aging-aware routing, it improves maximum age by about 33-34\% and age imbalance by 61-72\%, with the 3D case showing especially strong gains. However, both CAT and D-CAT rely on explicit aging tables and synchronizing aging information across the network, which can introduce memory overhead and limit scalability as the NoC size increases.

\section{Gaps in Prior Studies}

Prior studies on NoC routing algorithms has made significant progress in both fault-tolerant and aging-aware routing. However, there are still important gaps that this thesis aims to address. First, many aging-aware routing methods rely on explicit aging models or tables that require global state and synchronization, which can limit scalability and adaptability to changing conditions. Second, while RL-based fault-tolerant routing has shown promise, it has not been widely applied to aging-aware routing, which also requires dynamic adaptation to runtime conditions. The main trade-off across the literature is between model fidelity and implementation overhead. To this end, this work proposes to fill these gaps by developing a wear-aware routing algorithm that leverages linear approximate reinforcement learning to adaptively steer traffic away from stressed components without relying on explicit aging tables or models, while also maintaining acceptable network performance.